\subsection{Theme Synthesis — 生成、対話、行動、空間}

4月のmultimodalを通して見えるのは、画像・音声・視覚を一つの巨大モデルへ統合したという完成図ではなく、生成surface、full-duplex interaction、speech output、agent perception、spatial representationという要件が同時に立ち上がったことだ。7月のようなリアルタイム統合を語るにはまだ早い一方、その後の統合を必要にする部品はすでに揃い始めている。この段階差を残すことで、後続号の進展を単なるモデル名の更新ではなくsystem integrationの深化として読める。\autocite{src-0f697d747d9c,src-3fa7f7c1140f,src-d113b486a585,src-d33ba68aac3c,src-d8b7c1c36830,src-fec81b98fdc8}


\begin{claimboundary}
Claim Boundary: ChatGPT Images 2.0のreleaseとOpenAIが説明したscopeは一次資料で確認できるが、生成品質や性能優位は独立再現された評価ではない。\autocite{src-0f697d747d9c}
\end{claimboundary}

\begin{claimboundary}
Claim Boundary: MiniCPM-o 4.5、SpatialFusion、World2VLMなどの研究結果は著者報告であり、本稿は論文の問題設定と4月のchronologyを使う。比較結果を任意のmodel／hardwareへ一般化しない。\autocite{src-3fa7f7c1140f,src-d33ba68aac3c,src-fec81b98fdc8}
\end{claimboundary}

\begin{claimboundary}
Chronology Boundary: 後日改訂されたarXiv論文では、4月の初回投稿に存在した内容とlater revisionで追加・変更された本文を分離する。特にGLM-5V-Turboは保持locatorが2026年5月12日更新であり、その改訂内容を4月の出来事として遡及させない。\autocite{src-d113b486a585}
\end{claimboundary}

\begin{claimboundary}
Claim Boundary: Mistral newsの保存indexは具体的な4月entryを確認するために使い、現在のsite状態や不足するdetailから過去内容を推測しない。\autocite{src-d8b7c1c36830}
\end{claimboundary}
